\section{The clustered exponent and a range of inert row bounds}

% BEGIN I82B1
\paragraph{I82B1. Finite cohomology on a row with a uniform lift.}
At an inert row B point as in I82W1--I82W4 fix $w\mid p$, put
$d=[E_w:\mathbb Q_p]$, and use any coefficient field required in
the row test, with integers $O$, fraction field $L'$, and
$k_O=\operatorname{length}_O(O/pO)$. Such a field is good at
level one: $K=K_0$ and splitting the required cyclotomic level
implies splitting its level-one subextension. Write $F:T_1\to T_2$
for a tested lattice surjection. Its target has rank $r$, since
the tested representation is isomorphic to the specified twisted
dual of the rank-$r$ complement.

The module of finite cohomology classes $P=H^1_f(K,T_2)$ is
finitely generated over $O$ by the integral finiteness and
rationalization argument of A81R1. It needs at most $(d+1)r$
generators. Indeed, put $d_0=\dim_{L'} V_2^{G_K}$, with
$V_2=T_2[1/p]$. Its rational finite dimension is the dimension of
the de Rham tangent plus $d_0$ by Bloch--Kato, at most $dr+d_0$.
The invariants sequence for $T_2\subset V_2$ identifies the
torsion of $P$ with $(V_2/T_2)^{G_K}/(V_2^{G_K}/T_2^{G_K})$.
The denominator is a divisible module of corank $d_0$ and
splits off the numerator as an $O$-module, by injectivity over a
DVR. The numerator's uniformizer-torsion has dimension at most $r$
over the residue field, by inclusion in $V_2/T_2$. The finite
torsion quotient therefore needs at most $r-d_0$ generators.
This proves the assertion including its integral torsion.

If an integral finite lift exists for $p^e$ times every target
finite class, the row cokernel is killed by $p^e$ and has length
at most $(d+1)r\, e k_O$. Indeed the image contains $p^e P$.
On passing from classes to the source's cocycle test the target
integral coboundaries introduce no restriction: lift the defining
vector through the given onto lattice map and take its coboundary.
Thus the normalized bound on $\mathrm{RowAt}$ is $(d+1)r e$.

In I82W2's notation set $G=W_0+S_0+H$, $j=j_A$, and put
$h=B+H+S_0$ on an orbit with a thick letter, $h=G$ otherwise.
Every source and target jump is in $[-h,h]$, as are those of the
kernel by strictness. In the first case use $G\le2b'\le B$ for
the nonthick letters. The local proof of L81S2--L81S4, through
its estimate $e\le25h/p+3$, uses only that $K$ is unramified,
$p\ge5$, the maps are crystalline and onto, and these symmetric
jump bounds. It applies here without the split-prime declaration.
In particular the original local row inequality holds whenever
\[
  25h/p+3\le 3C_2 U,\qquad
 U=\frac{h_B}{(p-1)^{j+1}}+(G_P+1)(j+1)+L_{\rm large},
 \qquad G_P=w_P+S_c+\eta_3 .
\]
Indeed $d=n_\lambda$, and $\mathrm{coreB}=(d+1)rU$.
This core is positive and $X.\mathrm{lenRow}(i,w)$ is the
chosen rational strictly between $3C_2\mathrm{coreB}$ and
$(3C_2+1)\mathrm{coreB}$. This proves the full $\mathrm{RowAt}$
test for every orbit satisfying the displayed inequality.
In particular it holds on every orbit without a thick letter:
$h=G\le G_P$, $p\ge11$, $j\ge2$ and $C_2\ge1$.
% END I82B1

% BEGIN I82B2
\paragraph{I82B2. The numerical part of Theorem E at its original bound.}
For the general numerical theorem let $p\ge5$ be prime, $q=p-1$,
$j\ge1$, $g=\max(G',S',\eta',1)\le h$, with the integer parameters
of $\mathrm{ThmEBound}$. Let $\ell\ge1$ be the greatest arc
length. Both input height sequences $H_\nu$ have entries at most
$2h$ and at the first $j$ positions after each cut at most $2G'$.
Set $A_h=h/q^{j+1}$, $L_h=\log_p h$. For each sequence put
$c=\lceil1+\log_p(\max H_\nu+1)\rceil\le3+L_h$.
In the defining series at a cut, summation of the first $j$
terms and the geometric tail gives
\[
 \sum_{a\ge0}\frac{H_\nu(c_0+a+1)+c+3}{q^a}
 \le\frac q{q-1}(2G'+2h/q^j+c+3).
\]
In $r_F$ both the minimum and floors only lower the input to
multiplication by $p/q^2$. Hence its maximum $R$ over the
two sequences and cuts satisfies
\[
 R\le G'+4A_h+2L_h+6.
\]
For example $q/(q-1)\le4/3$, $p/q\le5/4$ and $p/q^2\le1/3$:
the coefficient of $G'$ is at most $8/9$, that of $A_h$ at most
$10/3$, and the remaining term is at most $(4/9)(L_h+6)+L_h+3$.

The source uses $M=\lfloor(h+2g)/q\rfloor+1$ and
$\sigma=\lceil(3g+M/p^{j-1})p/q^2+1\rceil$.
Since $p^{j-1}\ge q^{j-1}$ and $q^{j+1}\ge16$,
\[
 \sigma\le 2A_h+2g+3 .
\]
For the last polynomial valuation, each of the three intervals
is of length at most $3g$ in $[-2h,2h]$. At a level $p^k$ any
one interval has at most $3g/p^k+1$ nonzero multiples of $p^k$.
No such multiple occurs for $k>\log_p(2h)$. Summing in $k$
bounds the sum of valuations on an interval by $3g/q+\log_p(2h)$,
also if it crosses zero. The union has at most $9g$ elements.
By the definition of $\mathrm{vPP1}$, its supremum $\kappa$ costs
\[
 \kappa\le\ell+9g(\ell+1)+9g/q+3\log_p(2h)
          \le22g\ell+3L_h+3 .
\]
In particular the three separate interval positions cost three
logarithms independent of their separation.

With the original initial exponent $a_0=2\sigma+R+2g\ell$,
the full exponent is
\[
 e_{\rm par}=2R+4\sigma+2a_0+2\kappa
            =4R+8\sigma+4g\ell+2\kappa .
\]
Substitution and $1,G',g\le g\ell$ give
$e_{\rm par}\le32A_h+122g\ell+14L_h$.
Thus $\mathrm{ThmEBound}(800)$ holds with the original input
quantifiers (in fact $122$ suffices in its displayed upper bound).
The argument estimates the prescribed integer exponent before
asserting anything about the existence of a cocycle lift.
% END I82B2

% BEGIN I82B3
\paragraph{I82B3. Cuts and the initial inequality on every marked row B orbit.}
Suppose the orbit includes a thick letter. Use
$G'=G=W_0+H+S_0=g$, $h=B+H+S_0$ in I82B2. Keep the up/down
type and size on thick letters; assign type thin on all others.
The windows for a thin letter fit $[-G,G]$ by I82W2.
For the heights of both tested lattices one can take
$M_\nu(a)=M_a,\ h_\nu(a)=h_a$ from I82W2, since the $T_2$
jumps satisfy these same bounds. Their sums $H_\nu$ are at most
$2h$ everywhere and at most $2G$ off the marks.

Take all final marks of I82W4's clusters as cuts, including the
mark of a singleton. From one final mark to the first mark of
the next cluster there are at least $j$ intervening thin positions.
Between successive cuts the terminal cluster has length at most
$j$ from first mark to last. If the distance to the next cut is
greater than $3j+1$, insert a cut $j$ places after the current
one and repeat. After any such insertion the remaining distance
to the terminal mark exceeds $2j+1$, so the first mark of that
cluster is more than $j$ beyond the newly inserted cut. This
works also when the initial cyclic list has just one cut by
unfolding the full period once. The final gaps lie between $j$
and $3j+1$ inclusive. Give a cut its backward distance to the
preceding cut. All arcs are nonempty, lengths sum to $n_w$,
the $j$ positions after every cut are thin, and the backward arc
including its terminal cut but excluding its initial cut contains
at most the two marks of one cluster. These are $\mathrm{CutsOK}$,
$\mathrm{ThickLe2}$, and $\ell\le3j+1$.
Every thick size lies in $[B,B+H]$ and exceeds $g$ by I82W3;
this gives the required type-size bounds.

We can verify $\mathrm{WordOK}$ simultaneously for all the
allowed height sequences, at the same late family. I82B2 bounds
$R$ and $\sigma$ as above. Uniformly on this family,
\[
  \frac{h}{B}=1+o(1),\quad
  \frac{h}{Bq^{j+1}}=o(1),\quad
  \frac{g(3j+1)}B=o(1),\quad \frac{\log_p h}{B}=o(1).
\]
Indeed $G\le2b'$, $j\le L_w=O(\log x)$,
$b'/b=x\exp(-x^{\epsilon/3})$, and $H+S_0\le2b'$.
Also $b\to\infty$ from I82W1 and $b/b'\to\infty$, while
$p>x^{\epsilon_z/5}$ and $j\ge2$. In particular $2a_0<B$
eventually. For any two thick sizes $s_0,s_1$ in a backward
arc, $s_0\le B+H\le 2B$, $s_1\ge B$, and $p\ge11$. Thus
$s_0<p s_1$, $2s_0\le p s_1$, and $2a_0<s_1$.
These are both order-sensitive implications in $\mathrm{WordOK}$.
The enlarged cutoff is independent of $w$ and of the lattice map.

Consequently the only extra local assertion needed by the cut
method here is the actual integral finite lift
$F.\mathrm{LiftF}(p^{e_{\rm par}})$, writing the scalar explicitly,
on the marked orbits and maps just described. More precisely the
source writes $F.\mathrm{LiftF}(\mathrm{PD},e_{\rm par})$. If
that assertion holds, I82B1 and I82B2 bound the normalized
colength by $(d+1)r\cdot122(A_h+g\ell+L_h)$. To compare this
with the hub, $\log_p h\le h/q^{j+1}+j+1$: divide by $q^{j+1}$,
use $\log_p y\le y$ for $y>0$ and $\log_p q\le1$. Hence
$A_h+g\ell+L_h\le3U$, using $h\le h_B$, $g\le G_P$,
and $\ell\le3j+1$. Choosing the free common row constant
$C_2\ge122$ in F78D2 makes this at most $3C_2(d+1)rU$.
The resulting full row B colength on marked orbits outside
I82B1's range remains conditional on that integral lifting
assertion. I82B1 has already supplied the row there whenever
its displayed coarse-lift criterion holds.
% END I82B3

% BEGIN I82B4
\paragraph{I82B4. Integral truncation through the thin letters following a cut.}
Use the tame-invariant coordinate $v$ of L81S2, so that the
coefficient ring of each letter is $A=O[[v]]$,
$P_0=v+p$ and $\varphi(v)=vP_0^q$ times a unit, $q=p-1$.
Let $P_s=\varphi^s(P_0)$, which modulo $p$ is $v^{p^s}$
times a unit. Indeed substitution modulo $p$ sends $v$ to a
unit times $v^p$, and the coordinate has coefficients in
$\mathbb Z_p$. Up to units the $P_s$ are the primitive
cyclotomic factors in these tame coordinates. Put
\[
 M=\lfloor(h+2g)/q\rfloor+1,\quad
 f_j=v^M\prod_{s<j}P_s^{3g},\quad
 \sigma=\left\lceil(3g+M/p^{j-1})p/q^2+1\right\rceil .
\]
Let $N_c$ be the invariant Wach lattice of a tested crystalline
object at the cut $c$, and $P(c,j)=A\varphi^j N_c$ its span at
$c+j$. We first produce the exact truncation at this $\sigma$.
For any $z\in vN_c$ there exist $\tau\in P(c,j)$, an integer
$b\ge0$, and $x\in vN_c$ with
\[
  (\gamma-1)x-p^bz\in v^MN_c,\qquad
  p^b\tau-p^\sigma\varphi^j x\in f_j P(c,j).
\]
Select a free $A$-basis of $N_c$ and solve $(\gamma-1)y=z$
successively in degrees $1,\ldots,M-1$, with zero constant.
The $\gamma$ matrix is the identity modulo $v$ and the linear
term of its substitution is $u^q v$ for $u=1+p$. In degree
$k>0$ the divisor is $u^{qk}-1$, of valuation $1+v_p(k)$.
All earlier coefficients enter with integral multipliers.
Accordingly the degree-$k$ coefficient of this finite polynomial
primitive has denominator dividing $p^{w(k)}$,
$w(k)=k+v_p(k!)\le kp/q$.

View $p^\sigma\varphi^j(y)$ in the rational quotient
$P(c,j)/f_j P(c,j)$ with $p$ inverted. It has an integral
representative. It suffices to check one monomial $v^k$ and its
denominator. Write $D=3g+M/p^{j-1}$. If $qk\le D$, then
$w(k)\le Dp/q^2\le\sigma$. Otherwise
$\varphi^j(v)^k$ is a unit times
$v^k\prod_{s<j}P_s^{qk}$. Remove the factor $\prod_{s<j}P_s^{3g}$.
The remaining power $P_{j-1}^{qk-3g}$ has, modulo $v^M$,
$p$-valuation at least $qk-D$: expand
$P_{j-1}=pA_1+v^{p^{j-1}}A_2$ with $A_1,A_2$ integral.
Terms with more than $M/p^{j-1}$ copies of the second term
vanish. As $k$ increases past $D/q$, this gain grows with
slope $q$, whereas $w(k)\le kp/q$. Hence
$w(k)\le\sigma+\max(0,qk-D)$. This proves integrality in the
quotient term by term, also after multiplication by the integral
Frobenius-transformed basis. Choose the representative $\tau$,
take $b$ clearing all coefficients of $y$, and set $x=p^by$.
The quotient by $f_j$ has no $p$-torsion: $f_j$ is a unit times
a distinguished monic polynomial, and $P(c,j)$ is free on a
Frobenius-transformed basis. Thus the two asserted congruences
hold integrally.

Both $\gamma$ and its inverse preserve $(f_j)$, since they preserve
$(v)$ and $(P_0)$ and commute with $\varphi$. Also
$f_j\mid\varphi^j(v)^M$ since $qM>h+2g\ge3g$.
Commuting $\gamma$ with $\varphi$ in the preceding congruences,
then cancelling $p^b$ modulo $f_j$, gives
\[
 p^\sigma\varphi^j z-(\gamma-1)\tau\in f_jP(c,j)
                                      \subset v^M N_{c+j}.
\]
The last inclusion is integral: at the next $j$ thin letters
the lower height is at most $g$, so propagation of the
Frobenius span has at most $P_s^g$ as denominator at each step.
Finally $\varphi$ acts contractively on the product of the
$v^MN_i$. Relative to their bases each Frobenius step has the
extra factor $P_0^{qM-h}$ in the maximal ideal, and $\varphi$
preserves that ideal. The geometric series in this complete
module therefore inverts $1-\varphi$. These arguments apply to
the finite crystalline extension by the trivial representation as
well, since adjoining the zero jump preserves the given windows.
% END I82B4

% BEGIN I82B5
\paragraph{I82B5. Decomposing an actual finite target class into arc contributions.}
Write $d_\varphi E(i)=\varphi(E(i-1))-E(i)$ and
$d_\gamma E(i)=(\gamma-1)E(i)$. Represent a target finite class
by its integral extension of the trivial module by $T_2$.
The extension is crystalline, and R82W4 provides at every cut
$c$ an invariant Wach vector $x_c$ above $p^R$ times the
trivial generator, $R=rr_{\rm par}$. In the invariant \'etale
coordinates the resulting sequence of modules is exact by
B81L1 and tame averaging; identify its kernel with the target's
coordinates. Put $z_c=(\gamma-1)x_c\in vN_c(T_2)$.
Here saturation follows by writing $(\gamma-1)x_c=vx$ in the
extension lattice and using regularity of $v$ in the quotient,
then the Wach kernel identity A81R10. Apply I82B4 to get
$\tau_c\in P(c,j)$ and define, at the shifted cut $c+j$,
\[
 R_c=p^\sigma\varphi^j x_c-\tau_c,\qquad
 \xi_c=(\gamma-1)R_c\in v^MN_{c+j}(T_2).
\]
The projection of $R_c$ is $p^{R+\sigma}$ times the generator.
Extend $R_c$ forward by $\varphi$ from $c+j$ until just before
the next shifted cut, doing this around the whole cycle; call the
result $E(i)$. Thus $(\gamma-1)E$ lies everywhere in $v^MN_i(T_2)$.
For the preceding cut $c'=c-\ell_c$ put
\[
 w_c=\varphi^{\ell_c-j} R_{c'}-p^\sigma x_c,\qquad
 r_c=p^\sigma(\varphi^j w_c+\tau_c).
\]
The first difference projects to zero. Also $(\gamma-1)w_c$
is in $vN_c$, since its two contributions are a propagated
$\xi_{c'}$ and $-p^\sigma z_c$. These vectors have precisely the
backward arc denominator of $\mathrm{Wc}$. In detail, put
$\Lambda_c=\prod_{s<\ell_c}P_s^{g+h_{c-s}^{+}}$, where the
upper heights $h_i^{+}$ are $G$ at thin letters, $s_i+S_0$
at thick up letters and $H$ at thick down letters. Repeatedly
using $P_0^{h_i^{+}}N_i\subset A\varphi N_{i-1}$ gives
$\Lambda_c x_c\in P_Y(c',\ell_c)$ for the extension $Y$.
The propagated $R_{c'}$ is already in $P_Y(c',\ell_c)$.
The kernel in that span is $P_{T_2}(c',\ell_c)$: tensor the
kernel inclusion for Wach modules with $A$ via $\varphi^{\ell_c}$.
This base change is flat, since $A$ is finite free over
$\varphi(A)$ by distinguished division in the coordinate $v$
($\varphi(v)$ has residue order $p$); iterate. The spans after
base change embed in the \'etale module by invertibility of
their Frobenius matrices there. Hence
$\Lambda_c w_c\in P_{T_2}(c',\ell_c)$, as required.

The value $p^\sigma d_\varphi E$ is supported on the shifted cuts
and equals $r_c$ at $c+j$, by the definition of $w_c$.
Let $y_c$ be supported just there with that value. The sequence
$d_\gamma y_c$ lies in $\prod v^M N_i$, since at its single
support it equals
$p^\sigma d_\varphi d_\gamma E$ there. By I82B4 there is a
unique $z'_c$ in this product with $d_\varphi z'_c=d_\gamma y_c$.
Uniqueness gives $\sum_c z'_c=p^\sigma d_\gamma E$.
Each $(y_c,z'_c)$ is the source's $\mathrm{IsRho}$ pair for
$w_c$. Indeed its truncation for $(\gamma-1)w_c$ is
$-p^\sigma\tau_c$: the congruences of I82B4 scale by $-p^\sigma$,
and the extra propagated $\xi_{c'}$ in the first congruence
lies in $v^MN_c$.

Take a section $E_0(i)$ of the trivial generator in the invariant
\'etale extension. Its differences represent the given finite
class: subtract the coordinates of the original integral lift in
the Witt tensor and take its two differences, as in B81L3 and
A81R11. Such a section exists by exactness of the \'etale modules
and tame averaging. The sequence $p^{R+\sigma}E_0-E$ is
kernel-valued. Subtract the coboundary of $p^\sigma$ times this
sequence from $p^{R+2\sigma}(d_\varphi E_0,d_\gamma E_0)$.
What remains is exactly $\sum_c(y_c,z'_c)$.
This proves the native cut decomposition $\mathrm{Step1}$ on
the target extensions with its stated multiplier. The remaining comparison of source and target must lift these
individual arc data through $F$ modulo $vN_c$ at the prescribed
exponent. For the finiteness of the resulting source arc classes
see H82A1--H82A3.
% END I82B5
